
\section{Analysis}

\subsection{The Johnson-Lindenstrauss Guarantee}

\textbf{Figure 8.1} shows the distortion introduced by the random projection. 
For $r = 32$ and $d_h = 64$, the average inner-product distortion is:

\[
\frac{|\tilde{Q}_i \tilde{K}_j^\top - Q_i K_j^\top|}{|Q_i K_j^\top|} \approx 0.015
\]

\textbf{Observation.} At $r = 16$, the distortion reaches 32, 
which explains the accuracy drop.

\subsection{Why RPC Has Modest Efficiency Gains}

\textbf{Table 8.1: RPC's real efficiency at L=128.}

\begin{table}[H]
\centering
\small
\begin{tabular}{lll}
\toprule
\textbf{Component} & \textbf{Baseline} & \textbf{RPC (r=32)} \\
\midrule
Q/K/V proj & 0.5 ms & 0.5 ms \\
Sketch & 0 & 0.1 ms \\
Score & 0.4 ms & 0.25 ms \\
Softmax & 0.05 ms & 0.05 ms \\
@ V & 0.4 ms & 0.4 ms \\
Output proj & 0.2 ms & 0.2 ms \\
LoRA & 0 & 0.15 ms \\
\textbf{Total} & 1.55 ms & 1.65 ms \\
\bottomrule
\end{tabular}
\end{table}

RPC is actually \textbf{slower} at $L = 128$ because:
\begin{enumerate}
\item The sketch adds FLOPs.
\item LoRA adds FLOPs.
\item Kernel launch overhead dominates.

\end{enumerate}
\subsection{The Memory Problem}

\textbf{Critical.} RPC \textbf{does not reduce memory}. The score matrix 
$S^{RPC}$ is still $L \times L$. FlashAttention's memory savings come from 
not materializing $S$; RPC keeps this property but does not reduce the 
underlying $L^2$ structure.

\textbf{Figure 8.2} shows the memory usage vs sequence length. RPC and baseline 
overlap, while AASVD and EAT diverge.

\subsection{When Does RPC Help?}

RPC helps in the regime:

\[
L^2 d_h > L d_h r + L^2 r + \text{overhead}
\tag{8.1}
\]

Rearranging:

\[
L (d_h - r) > d_h r / L + \text{overhead} / L
\tag{8.2}
\]

For $d_h = 64$, $r = 32$, overhead is small at large $L$:

\[
L > \text{small} \Rightarrow L > 512
\]

So RPC helps at \textbf{long sequence lengths} ($L > 512$).

\subsection{Sketch Quality Analysis}

\textbf{Figure 8.3} shows the correlation between RPC scores and true scores as 
a function of $r$.

\begin{table}[H]
\centering
\small
\begin{tabular}{ll}
\toprule
\textbf{$r$} & \textbf{Correlation} \\
\midrule
8 & 0.84 \\
16 & 0.92 \\
32 & 0.96 \\
64 & 1.00 (identity) \\
\bottomrule
\end{tabular}
\end{table}

\textbf{Observation.} At $r = 32$, the correlation is 0.96. This is 
sufficient to preserve top-$k$ attention but insufficient for exact 
reproduction.

\subsection{Comparison with AASVD}

\begin{table}[H]
\centering
\small
\begin{tabular}{lll}
\toprule
\textbf{Aspect} & \textbf{AASVD} & \textbf{RPC} \\
\midrule
What is compressed & Weights & Q/K activations \\
Parameters & 50\% & +1\% \\
Memory & 40\% & 0\% \\
Attention FLOPs & Same & 50\% (at r=32) \\
Requires calibration & Yes & No \\
\bottomrule
\end{tabular}
\end{table}

\textbf{RPC and AASVD are complementary.} Combine both for parameter + FLOP 
reduction.

\subsection{Failure Mode Taxonomy}

\begin{table}[H]
\centering
\small
\begin{tabular}{lll}
\toprule
\textbf{Failure mode} & \textbf{Symptom} & \textbf{Root cause} \\
\midrule
Collapse & Acc $\approx$ 50.92\% & Weight transfer failure \\
Low correlation & Acc drops moderately & Sketch rank too small \\
Slow & Slower than baseline & Overhead dominates \\
No memory gain & VRAM unchanged & Score matrix still L $\times$ L \\
\bottomrule
\end{tabular}
\end{table}

\subsection{The `raise` Problem}

RPC raises \texttt{RuntimeError} in:
\begin{enumerate}
\item \texttt{copy\_weights} on missing projection.
\item \texttt{inject\_rpc} on any injection failure.

\end{enumerate}
This is \textbf{safer} than silent fallback but \textbf{less flexible} than A3/AASVD.

0.85

\subsection{Summary of Insights}

\begin{enumerate}
\item \textbf{RPC is not a compressor.}
\item \textbf{RPC does not reduce memory.}
\item \textbf{RPC helps only at long sequences (L > 512).}
\item \textbf{Theoretical speedup $\ne$ actual speedup.}
\item \textbf{\texttt{raise} prevents silent failure.}
\end{enumerate}
\begin{lstlisting}

---
\end{lstlisting}